\documentclass[10pt,a4paper]{amsart}
\usepackage[margin=19mm]{geometry}
\usepackage{amsmath,amssymb,mathtools}
\usepackage[scr=rsfs,cal=euler]{mathalfa}
\usepackage{xfrac}
\usepackage[unicode,hidelinks,bookmarks=false]{hyperref}
\newcommand{\ie}{i.e.\ }
\newcommand{\cf}{cf.\ }
\newcommand{\resp}{respectively\ }
\newcommand{\Subsection}{\subsection}
\providecommand{\nobreakdash}{\nobreak\hbox{-}}
\setlength{\emergencystretch}{2em}
\multlinegap0pt
\usepackage{bm}
\usepackage{bbm}
\usepackage{dsfont}
\usepackage{xspace}
\usepackage{cancel}
\usepackage{mathtools}
\usepackage{amsthm}
\makeatletter
\@ifundefined{theorem}{\newtheorem{theorem}{Theorem}}{}
\@ifundefined{lemma}{\newtheorem{lemma}[theorem]{Lemma}}{}
\@ifundefined{proposition}{\newtheorem{proposition}[theorem]{Proposition}}{}
\makeatother
\newcommand{\id}{{\rm id}}
\title[Invariant subalgebras rigidity for von Neumann algebras of g]{Invariant subalgebras rigidity for von Neumann algebras of groups arising as certain semidirect products}
\author{Tattwamasi Amrutam, Artem Dudko, Yongle Jiang, Adam Skalski}
\date{Source: arXiv:2507.12824v2 (CC BY 4.0)}
\begin{document}
\maketitle

\noindent\emph{Source and licence.} Invariant subalgebras rigidity for von Neumann algebras of groups arising as certain semidirect products. Version arXiv:2507.12824v2 (CC BY 4.0), distributed under the Creative Commons Attribution 4.0 International licence, as recorded in the arXiv licence metadata for this exact version and shown on the abstract page https://arxiv.org/abs/2507.12824v2. Full rights notice: licenses/arxiv-2507.12824-CC-BY-4.0.txt.\par\smallskip

\noindent\textit{Editorial scope.} Counted unit: the Proposition whose source label is \texttt{lemma: E(01 10) for cross product} in the article (source0.tex lines 730--732, source label \texttt{lemma: E(01 10) for cross product}), delivered together with the article's own complete original proof of it (source0.tex lines 733--799, one \texttt{proof} environment of 67 lines). No mathematics is written, repaired or re-derived: statement and proof are the article's own text line for line. Required context retained uncounted: prop: E properties (lines 233--242); prop:(GLinfty,F2)\_has\_ISR (lines 615--617); cor: E(v)=0 or v (lines 683--685); eq:E(01\_10) (lines 722--722). Every definition, display and label of those lines is retained unchanged. Omitted from this package and not redistributed: 1--232, 243--614, 618--682, 686--721, 723--729, 800--1199. Nothing inside a retained excerpt is omitted or altered: every retained body is delivered exactly as the article's lines are written, so no omission receipt of any kind accompanies this package. Citations: the retained excerpts carry no citation command at all, so no bibliography entry of the article is transcribed and no reference is redistributed. Reference adaptation: none was needed and none was made; every reference inside the delivered unit resolves inside it, so no unresolved reference or citation is produced. Printed slips kept literal: no printed word, symbol, hypothesis or number of the counted statement or of its proof was altered. Layout and class: the article's own class is \texttt{amsart} with packages including tikz-cd, babel, amsmath, amstext, amsfonts, amssymb, amsthm, amsrefs, mathtools, dsfont, color, graphicx, todonotes, float; the wrapper is the lane's pinned \texttt{amsart} editorial wrapper at 10pt on A4, which restates the theorem-like environments the retained text needs and adds the article's own macro definitions that the retained text needs (\texttt{\textbackslash id}), with extra wrapper packages (bm, bbm, dsfont, xspace, cancel, mathtools). The delivered numbering is the unit's own. Assets: no retained span contains a figure environment, a graphics-inclusion command, a PDF-page inclusion, a table or a TikZ picture; no image file is copied into this package, the package declares no asset dependency and no image file is redistributed with it. Licence: version arXiv:2507.12824v2 of this article is distributed under the Creative Commons Attribution 4.0 International licence, as recorded in the arXiv licence metadata for this exact version and shown on the abstract page https://arxiv.org/abs/2507.12824v2. A full rights notice is delivered at licenses/arxiv-2507.12824-CC-BY-4.0.txt. Characters: every retained excerpt is pure ASCII, so the delivered engine, XeLaTeX with its default Latin Modern text font, reports no missing character.
\begin{proposition}\label{prop: E properties}
	Let $\mathcal M\subseteq L(G)$ be a $G$-invariant von Neumann subalgebra, with the associated $\tau$-preserving conditional expectation $E: L(G)\rightarrow \mathcal M$. Then
	\begin{itemize}
		\item[(1)] $\tau(E(g)s)=\tau(E(g)E(s))=\tau(gE(s))$ for all $s, g\in G$;
		\item[(2)] $sE(g)s^{-1}=E(sgs^{-1})$ for all $s, g\in G$;
		\item[(3)] $sE(s^{-1})\in\mathcal M'\cap L(G)$ for any $s\in G$;
		\item[(4)] the formula $g \mapsto \tau(g^{-1}E(g))$ defines a character (normalised, $G$-invariant, positive-definite function) on $G$;
		\item[(5)] for any $g \in G$ we have $E(g)=0$ if and only if $\tau(g^{-1}E(g))=0$, and $E(g) = g$ if and only if $\tau(g^{-1}E(g))=1$.
	\end{itemize}
\end{proposition}

\begin{proposition}\label{prop:(GLinfty,F2)_has_ISR}
	The group $GL(\infty,F_2)$ has ISR.
\end{proposition}

\begin{lemma}\label{cor: E(v)=0 or v}
Either $E(u_v)=u_v$ for all $v\in F_2^\infty$, or $E(u_v)=0$ for all $v\in F_2^\infty\setminus \{0^\infty\}$.
\end{lemma}

\begin{equation}\label{eq:E(01_10)}E(u_s)=\sum\limits_{g\in GL(2, F_2)}u_g\cdot A_g,\;\;A_g\in L( F_2^2)=L^\infty(\widehat{ F}_2^2).\end{equation} 

\begin{proposition}\label{lemma: E(01 10) for cross product}
We have  $E(s)=s\cdot A$, where $A\in  L^\infty(\widehat{ F}_2^2)$ belongs to the set $\{0,1,[0,0]+[1,1]\}$.
\end{proposition}

\begin{proof}
The ideas in the first part of the proof, where we will establish that $E(s) = s \cdot A$ for some $A \in L^\infty(\widehat{ F}_2^2)$  are similar to these used in the proof of Proposition \ref{prop:(GLinfty,F2)_has_ISR}. 

Consider the expansion in the formula \eqref{eq:E(01_10)}. As in Proposition \ref{prop:(GLinfty,F2)_has_ISR}, let $\omega=(123)\in S_\infty\subseteq GL(\infty, F_2)$. Notice that $s^{-1} E(s)=sE(s)$ commutes with $E(\omega s\omega^{-1})=\omega E(s)\omega^{-1}=\omega sE(s)s^{-1}\omega^{-1}$. Consider the product
$$sE(s)\cdot \omega sE(s)s^{-1}\omega^{-1}=\sum\limits_{g\in GL(2, F_2)}sgA_g\cdot \sum\limits_{h\in GL(2, F_2)} \omega shA_hs^{-1}\omega^{-1}.$$
\vskip 0.1cm\noindent
${\bf 1)}$ Consider the element $g=\begin{bmatrix}
	1&0\\1&1
\end{bmatrix}$. Let $A_g=\sum\limits_{i,j\in\{0,1\}}c_{ij}\cdot [i,j]$, where $c_{ij}\in\mathbb C$. We have used the spanning property of $[i,j]$, i.e., $\{[w]: w\in \{0,1\}^k\}$ is a spanning set for $L(F_2^k)$. Taking into account that $C:=sg\cdot \omega sgs^{-1}\omega^{-1}$ has $C_{13}=1$, as  $\omega sgs^{-1}\omega^{-1} = \begin{bmatrix}
	1&0&0\\0&1&1\\0&0&1
\end{bmatrix}$ and so $C$ cannot appear with a nonzero coefficient in the expansion of $\omega sE(s)s^{-1}\omega^{-1}\cdot sE(s)$ we obtain:
$$sgA_g\cdot \omega sgA_gs^{-1}\omega^{-1}=0\;\;\Longrightarrow\;\; A_g\cdot h^{-1} A_gh=0,\;\;\text{where}\;\; h=(\omega sg)^{-1}=\begin{bmatrix}
	0&0&1\\0&1&1\\1&0&0
\end{bmatrix}.$$ 
Further, substituting the formula for $A_g$ we obtain:
\begin{align*}0&=A_g\cdot h^{-1} A_gh=\sum\limits_{i,j\in\{0,1\}}c_{ij}[i,j,\star]\cdot\sum\limits_{i,j\in\{0,1\}}c_{ij}[(i,j,\star)\cdot h]\\&=
	(c_{00}[0,0,\star]+c_{01}[0,1,\star]+c_{10}[1,0,\star]+c_{11}[1,1,\star])(c_{00}[\star,0,0]+c_{01}[\star,1,1]\\& \;\;\;\;\;+c_{10}[\star,0,1]+c_{11}[\star,1,0])\\&=
	c_{00}^2[0,0,0]+c_{01}^2[0,1,1]+c_{10}^2[1,0,1]+c_{11}^2[1,1,0]+
    \\& \;\;\;\;\;+c_{00}c_{10}[0,0,1]+c_{01}c_{11}[0,1,0]+c_{10}c_{00}[1,0,0]+c_{11}c_{01}[1,1,1].\end{align*}
   We conclude that $c_{ij}=0$ for all  $i,j\in\{0,1\}$, and so $A_g=0$.
\vskip 0.1cm\noindent
${\bf 2)}$  Consider now the element $t=\begin{bmatrix}
	0&1\\1&1
\end{bmatrix}$. Let $A_t=\sum\limits_{i,j\in\{0,1\}}c_{ij}\cdot [i,j]$, where $c_{ij}\in\mathbb C$. Observe that $sE(s)$ commutes with $\omega gE(s)g\omega^{-1}=E(\omega gsg^{-1}\omega^{-1})$, where $g=g^{-1}$ is the element from part $1)$. Similarly to case $1)$ we notice that for $C:=st\cdot \omega gtg^{-1}\omega^{-1}$ one has $C_{13}=1$ and so $C$ cannot appear with nonzero coefficient in the expansion of $\omega gE(s)g\omega^{-1}\cdot sE(s)$. We obtain:
$$stA_t\cdot \omega gtA_tg\omega^{-1}=0\;\;\Longrightarrow\;\; A_t\cdot h^{-1} A_th=0,\;\;\text{where}\;\; h=(\omega gt)^{-1}=\begin{bmatrix}
	0&0&1\\0&1&0\\1&0&0
\end{bmatrix}.$$ 
Similarly to case $1)$, substituting the formula for $A_h$ we obtain:
\begin{align*}0&=A_t\cdot h^{-1} A_th=\sum\limits_{i,j\in\{0,1\}}c_{ij}[i,j,\star]\cdot\sum\limits_{i,j\in\{0,1\}}c_{ij}[(i,j,\star)\cdot h]\\&=
	(c_{00}[0,0,\star]+c_{01}[0,1,\star]+c_{10}[1,0,\star]+c_{11}[1,1,\star])(c_{00}[\star,0,0]+c_{01}[\star,1,0]\\& \;\;\;\;\;+c_{10}[\star,0,1]+c_{11}[\star,1,1])\\&=
	c_{00}^2[0,0,0]+c_{01}^2[0,1,0]+c_{10}^2[1,0,1]+c_{11}^2[1,1,1]
    \\& \;\;\;\;\;+c_{00}c_{10}[0,0,1]+c_{01}c_{11}[0,1,1]+c_{10}c_{00}[1,0,0]+c_{11}c_{01}[1,1,0].\end{align*}
   % \textcolor{teal}{Artem says $\sum$-sign is missing. Answer: What do you mean by the sum? The terms are all explicitly written out. }
    We conclude that $c_{ij}=0$ for all  $i,j\in\{0,1\}$, and so $A_t=0$.

\smallskip
 Taking into account that $E(s)=sE(s)s^{-1}$, we obtain that $A_g=0$ for all $g\neq s$, which shows the statement claimed in the first line of the proof; we shall simply write $A$ for $A_s$ from now on.
\smallskip 
 
Notice that by the invariance of $E$ the element $s$ commutes with $A$, which means that $c_{01}=c_{10}$.
Further, by Proposition \ref{prop: E properties} (3), the element $A$ lies in $\mathcal M'$. Therefore, $A$ commutes with $\omega E(s)\omega^{-1}=\omega sA\omega^{-1}$. Let $A=\sum_{i,j\in\{0,1\}}c_{ij}[i,j]$. Then we have:
 \begin{align*}
 	A&\omega s A\omega^{-1}=\omega s A\omega^{-1}A\;\;\Longleftrightarrow\;\;\omega^{-1}A\omega\cdot A=A\cdot s\omega^{-1}A\omega s^{-1}\;\;&\Longleftrightarrow\\
 	\sum  &  c_{ij}([i,j]\omega)\cdot A=A\cdot\sum c_{ij}([i,j]\omega s^{-1})\;\;&\Longleftrightarrow\\ 
 	\sum &   c_{ij}[j,\star,i]\cdot \sum c_{kl}[k,l,\star]   =\sum c_{kl}[k,l,\star]\cdot \sum c_{ij}[\star,j,i]      \;\;&
 	\Longleftrightarrow \\ \sum & c_{ik}c_{kl}[k,l,i] =\sum c_{kl}c_{il}[k,l,i]   \;\;\Longleftrightarrow\;\; c_{ik}c_{kl}=c_{kl}c_{il}\;\;&\forall\;\;k,l,i\in\{0,1\}^3.\;\;\;\;\;\;\;\;\;\;\;\;\;\;
 \end{align*}
We will now analyse several cases.

\smallskip 
\noindent $a)$ If $c_{01}\neq 0$ (equivalently, $c_{10}\neq 0$), then taking $(k,l,i)=(0,1,0)$ and $(k,l,i)=(0,1,1)$ we obtain respectively that $c_{00}=c_{01}$ and $c_{10}=c_{11}$. Since all $c_{ij}$ are equal, one has $A=c_{00}\id$. Thus, $c_{00}s=E(s)=E(E(s))=c_{00}^2s$. It follows that $c_{00}\in \{0,1\}$ and $E(s)\in \{0,s\}$.

\smallskip\noindent
 $b)$ Suppose that $c_{01}=c_{10}=0$ but $c_{00}\neq 0$ and $c_{11}\neq 0$, so that $A \neq 0$. Since $E(s)(E(s))^*=(E(s))^2=A^2\in \mathcal M$, from Corollary \ref{cor: E(v)=0 or v} we obtain that $\mathcal M\supseteq L(F_2^\infty)$. Therefore, $s\cdot ([0,0]+[1,1])\in \mathcal M$. Now we may write $E(s)=sA=s(c_{00}[00]+c_{11}[11])\in \mathcal{M}$. Then $\mathcal{M}\ni E(s)\cdot[\star1]=s(c_{00}[01])$ since $[\star1]\in\mathcal{M}$. Similarly, $\mathcal{M}\ni E(s)\cdot [\star0]=sc_{11}[00]$. Then using $\langle s-sA, [01]\rangle=0=\langle s-sA,[00]\rangle$, we can deduce that $c_{00}=c_{11}=1$.

 
 \smallskip \noindent $c)$ Assume now that $E(s)=c_{00}s\cdot [0,0]$, where $c_{00}\neq 0$. Conjugating by the element $t=\begin{bmatrix}
 	0&1\\1&1
 \end{bmatrix}$ we obtain that $E(g)=c_{00}g\cdot [0,0]$, where $g=\begin{bmatrix}
 	1&0\\1&1
 \end{bmatrix}$. Let $g_k=I+e_{k1}$ for $k\geqslant 2$, so that $g_2=g$. Then $g_k$ commutes with $g$, and so with $E(g)$ as well. However, for $k\geqslant 3$, $$g_k\cdot [0,0]\cdot g_k= g_k\cdot ([0,0,\star^{k-3},0]+[0,0,\star^{k-3},1])\cdot g_k=[0,0,\star^{k-3},0]+[1,0,\star^{k-3},1]\neq [0,0].$$
 %\textcolor{teal}{Check the calculations above}
  We obtain a contradiction to the formula for $E(g)$, which shows that this case is impossible.

\smallskip \noindent $d)$ Finally, the case $E(s)=c_{11}s\cdot [1,1]$ can be treated similarly to the case $c)$.

\end{proof}	

\end{document}
